\uuid{2hM1}
\exo7id{12}
\titre{exo7 12}
\auteur{bodin}
\organisation{exo7}
\datecreate{1998-09-01}
\isIndication{false}
\isCorrection{false}
\chapitre{Nombres complexes}
\sousChapitre{Forme cartésienne, forme polaire}
\cpge{mpsi-mp2i.nombres-complexes.a, pcsi.nombres-complexes.a, ptsi.nombres-complexes.a, mpsi-mp2i.nombres-complexes.d, pcsi.nombres-complexes.d, ptsi.nombres-complexes.d, bcpst1.nombres-complexes.a, bcpst1.nombres-complexes.b}
\module{Algèbre}
\niveau{L1}
\difficulte{}

\contenu{
\question{
\'Ecrire sous la forme partie r\'eelle-partie imaginaire, puis sous la forme
module-argument le nombre complexe :
$$ \left(  \frac {1+i-\sqrt3(1-i)}{1+i}   \right)^2.$$
}
}
